% Options for packages loaded elsewhere
\PassOptionsToPackage{unicode}{hyperref}
\PassOptionsToPackage{hyphens}{url}
\documentclass[
]{article}
\usepackage{xcolor}
\usepackage{amsmath,amssymb}
\setcounter{secnumdepth}{-\maxdimen} % remove section numbering
\usepackage{iftex}
\ifPDFTeX
  \usepackage[T1]{fontenc}
  \usepackage[utf8]{inputenc}
  \usepackage{textcomp} % provide euro and other symbols
\else % if luatex or xetex
  \usepackage{unicode-math} % this also loads fontspec
  \defaultfontfeatures{Scale=MatchLowercase}
  \defaultfontfeatures[\rmfamily]{Ligatures=TeX,Scale=1}
\fi
\usepackage{lmodern}
\ifPDFTeX\else
  % xetex/luatex font selection
\fi
% Use upquote if available, for straight quotes in verbatim environments
\IfFileExists{upquote.sty}{\usepackage{upquote}}{}
\IfFileExists{microtype.sty}{% use microtype if available
  \usepackage[]{microtype}
  \UseMicrotypeSet[protrusion]{basicmath} % disable protrusion for tt fonts
}{}
\makeatletter
\@ifundefined{KOMAClassName}{% if non-KOMA class
  \IfFileExists{parskip.sty}{%
    \usepackage{parskip}
  }{% else
    \setlength{\parindent}{0pt}
    \setlength{\parskip}{6pt plus 2pt minus 1pt}}
}{% if KOMA class
  \KOMAoptions{parskip=half}}
\makeatother
\usepackage{longtable,booktabs,array}
\usepackage{caption}
\captionsetup[table]{skip=6pt}
\newcounter{none} % for unnumbered tables
\usepackage{calc} % for calculating minipage widths
% Correct order of tables after \paragraph or \subparagraph
\usepackage{etoolbox}
\makeatletter
\patchcmd\longtable{\par}{\if@noskipsec\mbox{}\fi\par}{}{}
\makeatother
\setlength{\emergencystretch}{3em} % prevent overfull lines
\providecommand{\tightlist}{%
  \setlength{\itemsep}{0pt}\setlength{\parskip}{0pt}}
\usepackage{bookmark}
\IfFileExists{xurl.sty}{\usepackage{xurl}}{} % add URL line breaks if available
\urlstyle{same}
% fallback for those not using the hyperref driver hyperxmp:
\makeatletter
\@ifundefined{xmpquote}{\newcommand{\xmpquote}[1]{#1}}{}
\makeatother
\hypersetup{
  pdftitle={Propuesta --- Proyecto 2},
  hidelinks,
  pdfcreator={LaTeX via pandoc}}

\title{Propuesta --- Proyecto 2}
\author{}
\date{}

\begin{document}
\maketitle

\textbf{Visualización reproducible en Python · Visualización I · Quinto
Semestre 2026}

\textbf{Equipo:} Equipo 7 --- Joshua Trujillo y Nicolás Figueroa
\textbar{} \textbf{Fecha:} 12 de octubre de 2026

\textbf{Docentes:} Adriana Pineda -- Paola Moreno \textbar{}
\textbf{Entrega del proyecto:} 26 de octubre de 2026

\section{1. Audiencia técnica
específica}\label{audiencia-tuxe9cnica-especuxedfica}

\textbf{Destinatario:} unidad de investigación comparativa sobre
determinantes del bienestar de una organización internacional (por
ejemplo, el grupo de investigación del PNUD o del Banco Mundial que
estudia indicadores de desarrollo y calidad de vida).

\textbf{Interés en la problemática:} este equipo necesita identificar de
manera fundamentada y reproducible qué combinaciones de condiciones
---PIB per cápita, apoyo social, esperanza de vida saludable, libertad,
generosidad y percepción de corrupción--- se asocian con altos niveles
de bienestar reportado, y en qué medida esas condiciones no alcanzan a
explicar la evaluación de vida de algunos países. El resultado orienta
futuras investigaciones sobre los determinantes explicados y no
explicados del bienestar y sirve como base verificable para estudios
comparados.

\section{2. Pregunta analítica}\label{pregunta-analuxedtica}

\begin{quote}
\textbf{¿Qué «recetas» de bienestar predominan entre los países con
mayor y menor evaluación de vida según el World Happiness Report 2024, y
en qué medida el componente no explicado (Distopía + residual) altera o
cuestiona esas recetas?}

Subpreguntas de apoyo:
\end{quote}

\begin{enumerate}
\def\labelenumi{\arabic{enumi}.}
\item
  ¿Qué combinación de los seis indicadores explicativos caracteriza a
  los países en los terciles superior e inferior del ladder score?
\item
  ¿Qué países presentan el componente Dystopia + residual más alto y qué
  tienen en común, sugiriendo que las variables disponibles no bastan
  para explicar su bienestar?
\item
  ¿Los países atípicos respecto de su «receta» esperada concentran altos
  o bajos niveles de residual, lo que indicaría que la explicación está
  incompleta?
\end{enumerate}

\section{3. Variables que se
emplearán}\label{variables-que-se-emplearuxe1n}

{\def\LTcaptype{none} % do not increment counter
\begin{longtable}[]{@{}
  >{\raggedright\arraybackslash}p{(\linewidth - 4\tabcolsep) * \real{0.3333}}
  >{\raggedright\arraybackslash}p{(\linewidth - 4\tabcolsep) * \real{0.3333}}
  >{\raggedright\arraybackslash}p{(\linewidth - 4\tabcolsep) * \real{0.3333}}@{}}
\toprule\noalign{}
\begin{minipage}[b]{\linewidth}\raggedright
\textbf{Variable (archivo oficial)}
\end{minipage} & \begin{minipage}[b]{\linewidth}\raggedright
\textbf{Significado}
\end{minipage} & \begin{minipage}[b]{\linewidth}\raggedright
\textbf{Uso previsto}
\end{minipage} \\
\midrule\noalign{}
\endhead
\bottomrule\noalign{}
\endlastfoot
Country name & País & Identificación y anotaciones en gráficos \\
Ladder score & Evaluación de vida (0--10) & Variable de interés / eje de
orden \\
upperwhisker / lowerwhisker & Intervalos del ladder score & Bandas de
incertidumbre (referencia) \\
Explained by: Log GDP per capita & Contribución del PIB per cápita &
Dimensión 1 de la «receta» \\
Explained by: Social support & Contribución del apoyo social & Dimensión
2 de la «receta» \\
Explained by: Healthy life expectancy & Contribución de la esperanza de
vida saludable & Dimensión 3 de la «receta» \\
Explained by: Freedom to make life choices & Contribución de la libertad
& Dimensión 4 de la «receta» \\
Explained by: Generosity & Contribución de la generosidad & Dimensión 5
de la «receta» \\
Explained by: Perceptions of corruption & Contribución de la percepción
de corrupción & Dimensión 6 de la «receta» \\
Dystopia + residual & Componente residual no explicado & Eje central de
la subpregunta 2 \\
\end{longtable}
}

\textbf{Criterio de selección:} se emplean las variables numéricas
disponibles en el archivo oficial WHR24\_Data\_Figure\_2.1.xls. No se
crearán indicadores derivados salvo la clasificación en terciles del
ladder score (pandas qcut, q=3), cuyo criterio se documentará en el
notebook. Los tres países con valores faltantes (Bahrain, Tajikistan y
State of Palestine) se conservarán y se excluirán explícitamente solo en
las visualizaciones que lo requieran.

\section{4. Plan de visualizaciones
(anteproyecto)}\label{plan-de-visualizaciones-anteproyecto}

{\def\LTcaptype{none} % do not increment counter
\begin{longtable}[]{@{}
  >{\raggedright\arraybackslash}p{(\linewidth - 4\tabcolsep) * \real{0.3333}}
  >{\raggedright\arraybackslash}p{(\linewidth - 4\tabcolsep) * \real{0.3333}}
  >{\raggedright\arraybackslash}p{(\linewidth - 4\tabcolsep) * \real{0.3333}}@{}}
\toprule\noalign{}
\begin{minipage}[b]{\linewidth}\raggedright
\textbf{\#}
\end{minipage} & \begin{minipage}[b]{\linewidth}\raggedright
\textbf{Visualización}
\end{minipage} & \begin{minipage}[b]{\linewidth}\raggedright
\textbf{Propósito analítico}
\end{minipage} \\
\midrule\noalign{}
\endhead
\bottomrule\noalign{}
\endlastfoot
Seaborn 1 & Distribución del ladder score y del residual & Forma global
y extremos \\
Seaborn 2 & Boxplot de los 6 indicadores por tercil de bienestar &
Comparación de «recetas» \\
Seaborn 3 & Heatmap: composición relativa (\%) de indicadores por grupo
& Perfil de receta por grupo \\
Seaborn 4 & Residual (Distopía + residual) vs. ladder score & Qué queda
sin explicar \\
Plotly 1 & Scatter interactivo: indicadores vs. ladder, tamaño =
residual & Explorar atípicos con hover \\
Plotly 2 & Barras top-15 con selector de indicador y color por tercil &
Profundizar por dimensión \\
\end{longtable}
}

\section{5. Coherencia esperada}\label{coherencia-esperada}

audiencia (unidad de investigación comparativa) ⇒ pregunta (recetas de
bienestar y componente no explicado) ⇒ variables (seis indicadores +
residual + ladder score) ⇒ preparación (renombre documentado, terciles,
manejo de faltantes) ⇒ visualizaciones (composición por grupos,
perfiles, residual, atípicos) ⇒ interpretación escrita tras cada gráfico
⇒ conclusiones (qué recetas predominan y qué queda sin explicar).

\textbf{Estado:} ☐ Borrador ☑ Enviado para aprobación ☐ Aprobado

\end{document}
